\documentclass[11pt]{article}
\usepackage{amsmath,amssymb,hyperref}
\hypersetup{pdfauthor={Daniel Alexander Trawin},pdfcreator={Trawin, Daniel Alexander},colorlinks=false}
\title{Phone2 Category Type Theory ClassicalRelay classicalBit}
\author{Daniel Alexander Trawin}
\date{2026-06-11T06:24:39Z}
\begin{document}
\maketitle
\section*{Metadata}
\begin{itemize}
\item TZPID: ID11051
\item UUID: 905eec64-26d5-5f10-afd8-0bbe3296d323
\item ORCID: \url{https://orcid.org/0009-0001-4630-3715}
\item Source: Quantum-Neural Cartographic Functor.txt line 63
\end{itemize}
\section*{Canonical Equation}
\[
"ClassicalRelay" → classicalBit,
\]
\section*{Statement}
Phone2 functor candidate in the category type theory family, extracted from Quantum-Neural Cartographic Functor.txt line 63, with canonical equation $"ClassicalRelay" → classicalBit,$.
\section*{Interpretation}
ID11051 records a Phone2 corpus equation in the category type theory family. The equation was atomized from Quantum-Neural Cartographic Functor.txt line 63. Nearest semantic context: (* Quantum-Classical Entanglement Functor *) Equation grounding: "ClassicalRelay" → classicalBit,
\end{document}
